This paper presents a corpus of 141,013 public Slovene-language posts by 432 authors on Bluesky, spanning August 2023 to mid-April 2026. The corpus was constructed through a protocol-native ATProto workflow combining relay-based discovery of Slovene-tagged posts, DID-resolved backfill of publicly available post histories from authors' Personal Data Servers, and post-level language filtering validated through manual annotation. As one of the first studies of Slovene-language use on Bluesky, the paper provides an initial descriptive account of an emerging public space on an open-protocol platform. The resulting corpus shows a marked increase in posting activity in late 2024. Reply posts account for 68.3\% of the material, with a high reply share already visible before this increase, suggesting that conversational use predated the surge. Further observed patterns include sparse and largely account-concentrated hashtag use, as well as content sharing shaped primarily by Slovene news media, GIF-based informal exchange, and Bluesky-native tooling. More broadly, the paper shows how protocol-level access in ATProto can support the construction of a high-precision corpus for an under-resourced language without relying on platform-specific research APIs. At the same time, it highlights legal and ethical questions that arise when small and highly identifiable language communities become reconstructable from public decentralised infrastructures.
